\section{A family of programs differing only in the linear term}
\label{sec:sweep}

An efficient frontier, a rolling rebalance and a scenario grid all solve the same
program repeatedly with a slightly different linear term: $G$, $C$, $b$ and
$m_{\mathrm{eq}}$ are fixed and only $a$ varies. Solved independently, each
instance rediscovers an active set it almost always already had.

What makes this exploitable is a reading of the invariants
\eqref{eq:inv1}--\eqref{eq:inv2}: \emph{neither factor depends on $a$}. $J$
depends on $G$ alone through $J\T G J = I$ and on the working set through the
orthogonal transformations accumulated into it; $R$ depends on $G$ and the working
set. The linear term enters only through $x_u = G^{-1}a$ and the right-hand sides.
So across such a family both factors are reusable verbatim, and the entire
solution can be recovered from them.

\subsection{Recovery}

\begin{proposition}[Recovery from cached factors]\label{prop:recover}
Let $(J,R)$ satisfy \eqref{eq:inv1}--\eqref{eq:inv2} for a working set $\Active$
of size $k$ with normals $A$. For a new linear term $a$, set
\begin{equation}
  x_u = J\,(J\T a),
  \qquad
  \rho = b_\Active - A\T x_u,
  \qquad
  R\T y = \rho,
  \qquad
  x = x_u + J_1 y,
  \qquad
  R\,u = y .
  \label{eq:recover}
\end{equation}
Then $(x,u)$ is the solution and multiplier pair of the working-set subproblem
\eqref{eq:sub}.
\end{proposition}

\begin{proof}
$x_u = JJ\T a = G^{-1}a$ by~\eqref{eq:inv1}. By Proposition~\ref{prop:blocks}(i),
$u = R^{-1}y = R^{-1}R^{-\top}\rho = (R\T R)^{-1}\rho = (A\T G^{-1}A)^{-1}\rho$,
which is the multiplier of~\eqref{eq:subsol}. For the iterate, $Ru = y$ gives
$J_1 y = J_1 R u = G^{-1}A u$ by Proposition~\ref{prop:blocks}(iii), so
$x = x_u + G^{-1}Au$, again as in~\eqref{eq:subsol}.
\end{proof}

The cost is $O(n^2)$, all of it in $x_u = J(J\T a)$: two matrix--vector products with
a dense $n \times n$ matrix, $J$ having lost its triangularity at the first update.
Everything after that is $O(nk + k^2)$. Re-deriving the factors for the
same working set instead costs $k$ insertions at $O(n^2)$ each, so the saving is a
factor of order $k$, largest exactly where it matters: a long-only optimum is a
vertex at which most bounds bind.

The $O(n^2)$ is worth insisting on, because the natural experiment cannot see it.
On a box family the active set grows with $n$, so $O(nk)$ and $O(n^2)$ predict the
same curve and either reading fits; holding $k$ fixed at $\qpSweepFixedK{}$ while
$n$ grows separates them, and the measured local exponent
$d\log t/d\log n$ then runs from $\qpSweepSlopeLo$ at $n = \qpSweepNlo{}$ to
$\qpSweepSlopeHi$ at $n = \qpSweepNhi{}$. Neither is $1$. The flat small-$n$ end is
an overhead floor, a hit there costing what its dozen array operations cost to
dispatch, and it ends where the $n^2$ overtakes dispatch.

Whether $(x,u)$ from~\eqref{eq:recover} is the answer to the \emph{full} program
is then one certificate away, and it is the same certificate as before: the
multipliers of the working-set inequalities must be non-negative and no constraint
outside the working set may be violated. By
Proposition~\ref{prop:sufficient} that is a proof. Note the asymmetry in cost: the
recovery is $O(n^2)$ but the check must look at every constraint, so it is
$O(nm)$ for a dense $C$, so verification and not recovery is what dominates a
successful reuse on a dense problem.

\subsection{Repair rather than restart}

When the check fails, the cached working set is stale, but it remains a far better
starting point than the unconstrained minimum. What blocks resuming from it is
\eqref{eq:invariant}: the iteration requires dual feasibility, and a stale set
typically has one or more negative multipliers. Those are precisely the
constraints that no longer belong.

\begin{proposition}[Repair terminates in a resumable state]\label{prop:repair}
Iterate the following from the cached set: recover $(x,u)$ by
\eqref{eq:recover}; if every inequality multiplier is non-negative, stop;
otherwise drop one constraint with a negative multiplier, downdate $(J,R)$ by
Section~\ref{ssec:delete}, and repeat. The loop stops after at most $k$ passes, in
a state satisfying~\eqref{eq:invariant}.
\end{proposition}

\begin{proof}
Each pass that does not stop shrinks the working set by one, so there are at most
$k$; with an empty set the sign condition is vacuous, which is the cold start of
Remark~\ref{rem:cold}. On exit $(x,u)$ solves the subproblem by
Proposition~\ref{prop:recover} and the sign conditions hold,
which is~\eqref{eq:invariant}; full column rank is inherited, deleting columns
preserving it.
\end{proof}

From such a state the iteration of Section~\ref{sec:iteration} cannot tell a
resumed solve from a cold one, the invariant being all it reads, so the resumed
walk needs only to drive the remaining primal infeasibility to zero. In the worst
case everything is dropped and the resumed walk is a cold solve, so the repair is
never worse than restarting up to the cost of the drops themselves.

\begin{remark}[Why this cannot change the answer]
Neither the reuse nor the repair introduces an approximation. The reuse returns a
point only when it carries a certificate, and the repair merely chooses a
different starting state for an iteration whose output is determined by
Proposition~\ref{prop:sufficient}. What does change is the reported iteration
count, which counts the working-set edits actually performed and is therefore not
comparable with a cold solve's.
\end{remark}
